\documentclass[11pt,a4paper]{article}

% ── Packages ──────────────────────────────────────────────────────────────────
\usepackage[utf8]{inputenc}
\usepackage[margin=1in]{geometry}
\usepackage{hyperref}
\usepackage{xcolor}
\usepackage{listings}
\usepackage{booktabs}
\usepackage{amsmath,amssymb}
\usepackage{fancyhdr}
\usepackage{enumitem}
\usepackage{mdframed}
\usepackage{tikz}

% ── Colours ───────────────────────────────────────────────────────────────────
\definecolor{codegreen}{rgb}{0,0.55,0}
\definecolor{codegray}{rgb}{0.5,0.5,0.5}
\definecolor{codepurple}{rgb}{0.3,0.5,0.8}
\definecolor{backcolour}{rgb}{0.96,0.96,0.96}
\definecolor{navy}{rgb}{0.0,0.2,0.55}
\definecolor{warnorange}{rgb}{0.85,0.4,0.0}
\definecolor{noteblue}{rgb}{0.1,0.35,0.65}

% ── Listing style ─────────────────────────────────────────────────────────────
\lstdefinestyle{mystyle}{
    backgroundcolor=\color{backcolour},
    commentstyle=\color{codegreen},
    keywordstyle=\color{navy}\bfseries,
    numberstyle=\tiny\color{codegray},
    stringstyle=\color{codepurple},
    basicstyle=\ttfamily\footnotesize,
    breaklines=true,
    captionpos=b,
    keepspaces=true,
    numbers=left,
    numbersep=5pt,
    showspaces=false,
    showstringspaces=false,
    tabsize=2
}
\lstset{style=mystyle}

% ── Warning box ───────────────────────────────────────────────────────────────
\newmdenv[
  linecolor=warnorange,
  backgroundcolor=warnorange!10,
  linewidth=2pt,
  roundcorner=5pt
]{warningbox}

% ── Note box ──────────────────────────────────────────────────────────────────
\newmdenv[
  linecolor=noteblue,
  backgroundcolor=noteblue!8,
  linewidth=2pt,
  roundcorner=5pt
]{notebox}

% ── Header / Footer ───────────────────────────────────────────────────────────
\pagestyle{fancy}
\fancyhf{}
\rhead{\textbf{Crazyflie 2.1 — Session 2 Documentation}}
\lhead{July 28, 2026}
\rfoot{Page \thepage}

% ── Title ─────────────────────────────────────────────────────────────────────
\title{\textbf{\Huge Crazyflie 2.1 ROS 2 Humble}\\[0.4em]
\Large Session 2: Thrust Control, Recovery State Diagnosis,\\
and Closed-Loop Anti-Drift P-Controller}
\author{\textbf{Karan Kumar Anand}\\
Department of Computer Science \& Engineering\\
IIIT Delhi}
\date{July 28, 2026}

\begin{document}
\maketitle

\begin{abstract}
This document covers the second development session for the Crazyflie 2.1 autonomous flight project. The primary findings of the session include the diagnosis and resolution of the Kalman filter \emph{recovery state} blocking motor output on a zero-deck drone, the development of a raw vertical thrust node (\texttt{vertical\_thrust.py}), and the implementation of a closed-loop proportional (P) position controller (\texttt{anti\_drift\_controller.py}) that subscribes to the drone's IMU-estimated pose and continuously corrects altitude, pitch, and roll to resist positional drift.
\end{abstract}

\tableofcontents
\newpage

% ─────────────────────────────────────────────────────────────────────────────
\section{Session Summary and Accomplishments}
% ─────────────────────────────────────────────────────────────────────────────

\begin{table}[h!]
\centering
\caption{Work completed on July 28, 2026}
\vspace{0.3em}
\begin{tabular}{@{}lll@{}}
\toprule
\textbf{\#} & \textbf{Item}                           & \textbf{Status} \\ \midrule
1           & Diagnosed Crazyradio USB ``Resource busy'' error  & ✓ Resolved \\
2           & Diagnosed \emph{recovery state} blocking motor output & ✓ Resolved \\
3           & Confirmed \texttt{/cf231/pose} published no data without Kalman reset & ✓ Identified \\
4           & Implemented \texttt{vertical\_thrust.py} with Kalman reset      & ✓ Done \\
5           & Implemented \texttt{anti\_drift\_controller.py} (P-controller)  & ✓ Done \\
6           & Created \texttt{QUICKSTART.txt} cheat sheet                     & ✓ Done \\
7           & Registered all new nodes in \texttt{setup.py}                  & ✓ Done \\
\bottomrule
\end{tabular}
\end{table}

% ─────────────────────────────────────────────────────────────────────────────
\section{Problem Diagnosis: Why the Drone Would Not Fly}
% ─────────────────────────────────────────────────────────────────────────────

\subsection{Crazyradio USB Resource Busy}

When attempting to start the ROS 2 server after a previous run was interrupted with \texttt{Ctrl+C}, the following error appeared:

\begin{lstlisting}[language=bash]
[crazyflie_server-2] what(): No Crazyradio dongle found!
ERROR: cflib.crazyflie: Couldn't load link driver: [Errno 16] Resource busy
\end{lstlisting}

\textbf{Root cause:} The Linux kernel's USB driver retains a lock on the Crazyradio device after a crash or forced process kill. The fix was:

\begin{lstlisting}[language=bash]
# Force-kill any leftover process holding the USB port
sudo pkill -9 -f crazyflie_server

# Unplug the Crazyradio dongle physically, wait 3 seconds, replug
# Then relaunch the server
ros2 launch crazyflie launch.py
\end{lstlisting}

\subsection{The Kalman Filter Recovery State (Critical Finding)}

Even after the server started successfully and all ROS 2 services were available (including \texttt{/cf231/arm} and \texttt{/cf231/takeoff}), the drone's motors did not respond. Investigation revealed:

\begin{lstlisting}[language=bash]
# /cf231/pose produced NO data output:
ros2 topic echo /cf231/pose
# <--- completely silent, no messages
\end{lstlisting}

\textbf{Root cause explained:}

The Crazyflie 2.1 runs an onboard Extended Kalman Filter (EKF) to estimate its position. This EKF requires at least one of the following position inputs to converge:
\begin{itemize}
    \item Flow Deck v2 (optical flow + distance sensor)
    \item Lighthouse positioning system
    \item Loco Positioning System (UWB)
    \item External motion capture system (Vicon/Qualisys)
\end{itemize}

Since this drone has \textbf{zero expansion decks}, the EKF has no sensor input. Without convergence, the firmware enters a \textbf{recovery/safety state} and blocks all motor output — even when thrust is explicitly commanded via \texttt{cmd\_vel\_legacy}.

\begin{notebox}
\textbf{Key Insight:} The high-level commander (\texttt{takeoff}, \texttt{land}, \texttt{sequence}) requires a fully converged Kalman filter. The low-level legacy commander (\texttt{cmd\_vel\_legacy}) also requires the recovery state to be cleared via a manual Kalman filter reset before it will accept thrust commands.
\end{notebox}

% ─────────────────────────────────────────────────────────────────────────────
\section{Solution: Kalman Filter Reset via ROS 2 Parameter Service}
% ─────────────────────────────────────────────────────────────────────────────

The recovery state is cleared by writing firmware parameter \texttt{kalman.resetEstimation = 1} through the \texttt{/crazyflie\_server/set\_parameters} ROS 2 service. This tells the firmware:

\begin{quote}
\emph{``Assume you are currently at position $(0, 0, 0)$ on the ground. Accept motor commands from this point forward.''}
\end{quote}

\begin{lstlisting}[language=python]
from rcl_interfaces.srv import SetParameters
from rcl_interfaces.msg import Parameter, ParameterValue, ParameterType

def reset_kalman(self):
    param = Parameter()
    param.name = "cf231.params.kalman.resetEstimation"
    param.value = ParameterValue(
        type=ParameterType.PARAMETER_INTEGER,
        integer_value=1
    )
    req = SetParameters.Request()
    req.parameters = [param]
    future = self.param_client.call_async(req)
    rclpy.spin_until_future_complete(self, future)
    time.sleep(2.0)  # allow filter to re-initialise at (0,0,0)
\end{lstlisting}

After this call, the \texttt{/cf231/pose} topic begins publishing Kalman-estimated position based on IMU dead-reckoning.

\begin{warningbox}
\textbf{Limitation:} Without a physical positioning deck, the Kalman filter relies purely on IMU dead-reckoning (integrating accelerometer data). This estimate drifts significantly after $\sim$10--15 seconds. A Flow Deck v2 is required for reliable long-duration hover.
\end{warningbox}

% ─────────────────────────────────────────────────────────────────────────────
\section{New Script 1: \texttt{vertical\_thrust.py} — Open-Loop Raw Thrust}
% ─────────────────────────────────────────────────────────────────────────────

\subsection{Purpose}
This node applies a fixed raw thrust value through \texttt{cmd\_vel\_legacy} to test whether the drone can physically lift off. It is \emph{open-loop} — no position feedback is used, so the drone will drift freely.

\subsection{Execution Sequence}
\begin{enumerate}
    \item Reset the Kalman filter (exit recovery state).
    \item Arm the drone via \texttt{/cf231/arm}.
    \item Publish raw thrust at a configurable value for a configurable duration.
    \item Thrust returns to 0.
\end{enumerate}

\subsection{Twist Message Mapping (cmd\_vel\_legacy)}

\begin{table}[h!]
\centering
\caption{Mapping of \texttt{geometry\_msgs/Twist} fields to drone commands}
\vspace{0.3em}
\begin{tabular}{@{}lll@{}}
\toprule
\textbf{Twist Field} & \textbf{Drone Axis} & \textbf{Units} \\ \midrule
\texttt{linear.x}    & Pitch (forward/backward) & degrees \\
\texttt{linear.y}    & Roll (left/right)        & degrees \\
\texttt{linear.z}    & Thrust                   & 0 -- 65535 \\
\texttt{angular.z}   & Yaw rate                 & degrees/s \\ \bottomrule
\end{tabular}
\end{table}

\subsection{Thrust Calibration Table}

\begin{table}[h!]
\centering
\caption{Thrust values and expected physical behaviour}
\vspace{0.3em}
\begin{tabular}{@{}cll@{}}
\toprule
\textbf{Thrust Value} & \textbf{Battery Level} & \textbf{Expected Result} \\ \midrule
38\,000 & Any   & Motors spin, no lift-off \\
42\,000 & 100\% & Approximate hover \\
46\,000 & 50\%  & Approximate hover \\
50\,000 & Low   & Strong lift-off (caution) \\
$>55\,000$ & --- & Dangerous, do not use indoors \\ \bottomrule
\end{tabular}
\end{table}

\subsection{Run Command}
\begin{lstlisting}[language=bash]
ros2 run my_crazyflie vertical_thrust
\end{lstlisting}

% ─────────────────────────────────────────────────────────────────────────────
\section{New Script 2: \texttt{anti\_drift\_controller.py} — Closed-Loop P-Controller}
% ─────────────────────────────────────────────────────────────────────────────

\subsection{Motivation}

When using only vertical thrust (open-loop), the drone drifts in the horizontal plane because small asymmetries in motor speeds, propeller balance, and centre-of-gravity cause uneven forces. With no feedback, there is no mechanism to detect or compensate this drift.

A \textbf{proportional (P) controller} uses real-time position feedback to generate corrective roll and pitch commands, and adjusts thrust based on altitude error.

\subsection{Control Law}

Let the desired target position be $(x_d, y_d, z_d)$ and the current Kalman-estimated position be $(x, y, z)$. The position errors are:

\begin{align}
e_x &= x_d - x \quad \text{(forward error)}\\
e_y &= y_d - y \quad \text{(lateral error)}\\
e_z &= z_d - z \quad \text{(altitude error)}
\end{align}

The P-controller outputs:

\begin{align}
\text{thrust} &= T_{\text{hover}} + K_{p,z} \cdot e_z \\
\text{pitch}  &= K_{p,xy} \cdot e_x \\
\text{roll}   &= -K_{p,xy} \cdot e_y
\end{align}

where $T_{\text{hover}} = 42000$ is the baseline hover thrust, $K_{p,z} = 8000\ \text{(thrust units/m)}$ and $K_{p,xy} = 5.0\ \text{(degrees/m)}$.

\subsection{Controller Gains}

\begin{table}[h!]
\centering
\caption{Tunable controller parameters}
\vspace{0.3em}
\begin{tabular}{@{}llll@{}}
\toprule
\textbf{Parameter} & \textbf{Default} & \textbf{Effect if increased} & \textbf{Effect if decreased} \\ \midrule
\texttt{HOVER\_THRUST}  & 42\,000     & More aggressive lift    & Drone sinks \\
\texttt{Kp\_z}          & 8\,000      & Faster altitude correction & Slow altitude response \\
\texttt{Kp\_xy}         & 5.0         & Stronger drift correction  & Slower drift correction \\
\texttt{TARGET\_Z}      & 0.4\,m      & Higher hover altitude  & Lower hover altitude \\ \bottomrule
\end{tabular}
\end{table}

\subsection{ROS 2 Data Flow}

The closed-loop control system consists of three components communicating over ROS 2 topics and services:

\begin{center}
\renewcommand{\arraystretch}{1.6}
\begin{tabular}{ccc}
\fbox{\begin{minipage}{3.8cm}\centering\textbf{anti\_drift\_controller}\\{\small (Terminal 2)}\end{minipage}}
& $\xrightarrow{\quad\texttt{/cf231/cmd\_vel\_legacy}\quad}$
& \fbox{\begin{minipage}{3.8cm}\centering\textbf{crazyflie\_server}\\{\small (Terminal 1)}\end{minipage}} \\
& &  $\downarrow$ {\small Crazyradio USB}  \\
\raisebox{1.2em}{\small $\xleftarrow{\quad\texttt{/cf231/pose}\text{ (Kalman estimate)}\quad}$}
& &
\fbox{\begin{minipage}{3.8cm}\centering\textbf{Crazyflie 2.1}\\{\small (hardware)}\end{minipage}}
\end{tabular}
\end{center}

\vspace{0.5em}
\noindent The controller publishes corrective roll, pitch, and thrust values on \texttt{/cf231/cmd\_vel\_legacy} at 50\,Hz. The drone's IMU-based Kalman filter feeds back its estimated position on \texttt{/cf231/pose}, which the controller subscribes to in order to compute position errors.

\subsection{Flight Phases}

\begin{enumerate}
    \item \textbf{Phase 1 — Ramp-up ($\sim$1 s):} Thrust increments from 0 to \texttt{HOVER\_THRUST} in steps of 2000 to avoid aggressive jolts.
    \item \textbf{Phase 2 — P-control loop (10 s, 50 Hz):} Position errors are computed from \texttt{/cf231/pose}. Roll, pitch, and thrust corrections are calculated and published every 20\,ms.
    \item \textbf{Phase 3 — Ramp-down:} Thrust decrements from \texttt{HOVER\_THRUST} to 0 for a soft landing.
\end{enumerate}

\subsection{Run Command}
\begin{lstlisting}[language=bash]
ros2 run my_crazyflie anti_drift_controller
\end{lstlisting}

% ─────────────────────────────────────────────────────────────────────────────
\section{Crazyradio Hardware Exclusivity Rule}
% ─────────────────────────────────────────────────────────────────────────────

A key operational constraint identified during this session is that the Crazyradio USB dongle is a \textbf{single shared hardware resource}. Only one program may hold the radio connection at a time.

\begin{table}[h!]
\centering
\caption{Mutually exclusive connection modes}
\vspace{0.3em}
\begin{tabular}{@{}lcc@{}}
\toprule
\textbf{Program running}        & \textbf{ROS 2 nodes work?} & \textbf{cfclient works?} \\ \midrule
\texttt{ros2 launch crazyflie launch.py} & \checkmark Yes           & $\times$ No \\
\texttt{cfclient} (GUI)                  & $\times$ No              & \checkmark Yes \\
\texttt{battery\_check.py}               & $\times$ No              & $\times$ No \\
Nothing running                          & ---                      & --- \\ \bottomrule
\end{tabular}
\end{table}

% ─────────────────────────────────────────────────────────────────────────────
\section{Multi-Terminal Operation Pattern}
% ─────────────────────────────────────────────────────────────────────────────

Multiple ROS 2 nodes may run simultaneously, but only \textbf{one flight controller} must be active at any time to avoid conflicting commands on \texttt{/cf231/cmd\_vel\_legacy}.

\begin{lstlisting}[language=bash]
# Terminal 1 — always running
ros2 launch crazyflie launch.py

# Terminal 2 — active flight controller (only one at a time)
ros2 run my_crazyflie anti_drift_controller

# Terminal 3 — safe to run in parallel (read-only monitoring)
ros2 topic echo /cf231/pose

# Terminal 4 — emergency stop at any time
ros2 run my_crazyflie emergency
\end{lstlisting}

% ─────────────────────────────────────────────────────────────────────────────
\section{Complete Node Inventory (\texttt{my\_crazyflie} Package)}
% ─────────────────────────────────────────────────────────────────────────────

\begin{table}[h!]
\centering
\caption{All registered ROS 2 nodes in \texttt{my\_crazyflie}}
\vspace{0.3em}
\begin{tabular}{@{}lll@{}}
\toprule
\textbf{Command}                                 & \textbf{Script}                   & \textbf{Description} \\ \midrule
\texttt{ros2 run my\_crazyflie arm}              & \texttt{arm.py}                   & Arm motors only \\
\texttt{ros2 run my\_crazyflie takeoff}          & \texttt{takeoff.py}               & Arm + takeoff to 0.3\,m \\
\texttt{ros2 run my\_crazyflie land}             & \texttt{land.py}                  & Safe landing \\
\texttt{ros2 run my\_crazyflie emergency}        & \texttt{emergency.py}             & Instant motor kill \\
\texttt{ros2 run my\_crazyflie sequence}         & \texttt{sequence.py}              & Full auto flight loop \\
\texttt{ros2 run my\_crazyflie vertical\_thrust} & \texttt{vertical\_thrust.py}      & Raw open-loop thrust \\
\texttt{ros2 run my\_crazyflie anti\_drift\_controller} & \texttt{anti\_drift\_controller.py} & Closed-loop P hover \\
\texttt{ros2 run my\_crazyflie go\_to}           & \texttt{go\_to.py}                & Move to position \\
\texttt{ros2 run my\_crazyflie velocity\_control}& \texttt{velocity\_control.py}     & Manual thrust/velocity \\
\bottomrule
\end{tabular}
\end{table}

% ─────────────────────────────────────────────────────────────────────────────
\section{Next Steps and Recommendations}
% ─────────────────────────────────────────────────────────────────────────────

\begin{enumerate}
    \item \textbf{Tune \texttt{HOVER\_THRUST}:} Increase from 42\,000 toward 46\,000 if the drone cannot lift off, or decrease if it climbs too aggressively.
    \item \textbf{Tune P gains:} Increase \texttt{Kp\_z} if altitude correction is sluggish. Increase \texttt{Kp\_xy} if sideways drift is not corrected quickly enough.
    \item \textbf{Install Flow Deck v2:} For reliable, long-duration hover without IMU drift, a Flow Deck v2 is strongly recommended. With it, \texttt{sequence.py} and \texttt{takeoff.py} will also function correctly.
    \item \textbf{Upgrade to PID controller:} After the P-controller is verified working, add integral (I) and derivative (D) terms to eliminate steady-state error and overshoot.
    \item \textbf{Trajectory following:} After stable hover is confirmed, implement waypoint navigation using a sequence of \texttt{TARGET\_X/Y/Z} setpoints.
\end{enumerate}

% ─────────────────────────────────────────────────────────────────────────────
\section{Conclusion}
% ─────────────────────────────────────────────────────────────────────────────

This session resolved the critical blocker preventing the Crazyflie 2.1 from flying in the zero-deck configuration: the Kalman filter recovery state. By programmatically resetting the onboard estimator via the ROS 2 parameter service, motor commands are now accepted. A raw thrust node and a closed-loop proportional position controller have been implemented and documented, providing a solid foundation for further research into autonomous flight control.

\end{document}
